% !TeX root = main.tex

\chapter{Low Dimensions}

\section{Two-dimensional Points ($d=2$)}

Numerous works have explored the two-dimensional case, and there exist well-known
$\mathcal{O}(n \log n)$ solutions. Perhaps the most famous one is through divide-and-conquer,
so let's begin with that.

\subsection{Divide-and-conquer}

Let's divide the set of points into two (approximately equal) halves. If $n$ is odd, the
first set will just be of size $\lceil \frac{n}{2} \rceil$ and the second set will
be of size $\lfloor \frac{n}{2} \rfloor$. Intuitively, we want those with smaller $x$
in one set and those with larger $x$ in the other. For convenience, we can sort the
points beforehand, with the sorting key being the point vector itself, i.e., perform element-wise
comparison (check who has the smaller $x$, else if they have equal $x$, compare by $y$ instead).

In the base case of $n = 1$ we simply return infinity (there's no pair!), else let's assume we have solutions to
both halves. Let's take the minimum answer from each half (we'll denote this distance $\delta$).
Our goal is to potentially connect points from the two sets to achieve a new minimum. Consider the
vertical divider ($x = x_{\text{mid}}$) used to partition the halves. Our candidate points are going to lie in the interval
$(x_{\text{mid}} - \delta, x_{\text{mid}} + \delta)$. For some point $(x_{0}, y_{0})$ in this set, we may check just those in the
interval $y \in (y_{0}-\delta, y_{0} + \delta)$. In fact we may just check those points with $y \leq y_{0}$,
avoiding duplicate work.

This ``box'' of width $2 \cdot \delta$ and height $\delta$ contains $\mathcal{O}(1)$ points. Consider
the left and right squares composing this box. If we partitioned each square into a grid of four smaller
sub-squares of width $\frac{\delta}{2}$, you may fit at most one point into each sub-square. Even
if we had placed points as far apart as possible, the distance is still $\frac{\delta}{\sqrt{2}}$.
If two points fitted in a sub-square, that'll imply we should have gotten a smaller answer
from one of the sub-problems, producing a contradiction.

To check points in this box easily, we'd like the points to be sorted by the $y$ coordinate instead.
This can be done exactly like how you'd perform \textit{Merge Sort}. Surprisingly, the nature
of our divide-and-conquer approach is no different from that of Merge Sort, and our time complexity
is hence $\mathcal{O}(n \log n)$.

\subsection{Sweepline}

One downside of divide-and-conquer is the runtime overhead of the recursive calls. An alternative
that yields the same time complexity (and also happens to be a bit faster empirically) is
a sweepline approach. We begin by sorting the points like in divide-and-conquer. Next, we'll
iterate through the points, and maintain a set of earlier points. This set organises the points
in terms of $y$ values.

Given the current minimum distance $\delta$ and the current
point $(x_{0}, y_{0})$, we only need to check points in the set with $y \in (y_{0} - \delta, y_{0} + \delta)$.
Using a similar box analogy, one may prove we check a constant number of points, producing an
overall complexity of $\mathcal{O}(n \log n)$. Implementation-wise, we use C\texttt{++}'s
\texttt{std::set} that uses a balanced binary search tree (specifically a red-black tree) under the hood.

\subsection{Grid Decomposition}

As alluded to at the start, there are algorithms that offer linear expected time. The trick is
to visualise the search space as a grid, and group points into squares on a grid. To figure
out the nearest pair, there are two cases to handle. First, iterate all pairs within
a grid square. Second, considering a point within a square, let's check its distance
against all other points in adjacent squares. If we pick a good grid width, we'll only
run through a constant number of pair checks per grid square, thus yielding $\mathcal{O}(n)$ time.

Intuitively, the grid width should be neither too big nor too small. A grid that's too fine-grained
might make it impossible to find a neighbouring point in the vicinity, and a large grid width
just degenerates into the brute force search. In a perfect world, we'd take the distance between
the closest pair as the grid width, but that brings us right back to the original problem. Eventually,
people discovered that taking the minimum distance of a bunch of random points (i.e., sample $n$ pairs)
produces a grid width that works surprisingly well.

Implementation-wise, we use the \textbf{SplitMix64} hash function paired with a hash map. At
the moment we settled with the standard library's \code{std::unordered_map}. Despite
the amortised $\mathcal{O}(1)$ lookup, its closed addressing scheme isn't particularly
cache-friendly.

\section{Beyond Two Dimensions}

While the two-dimensional algorithms have proven encouraging, both theoretically and empirically,
the real challenge is to solve the problem at higher dimensions, and this is where tools like
divide-and-conquer begin to fall apart as you can't easily generalise them.

\subsection{Grid Decomposition (Revisited)}

One neat property of grid decomposition is that it extends naturally to arbitrary dimensions.
A square morphs into a hypercube, and now we need to check adjacent hypercubes instead. There's
one big problem though; at dimension $d$, there are $\frac{3^{d} - 1}{2}$ neighbours to check, and
beyond a certain point, you're better off just brute-forcing.

\subsection{$k$-dimensional Tree}

A closely related problem to finding the closest pair is to handle point queries where given a
certain point, we ask for the nearest point to it. Answering the main problem then becomes
calling point queries on every point, and if we can perform this point query fast enough, we'll
beat the brute force approach.

As fate would have it, decades ago people already invented a data structure called the $k$-dimensional
tree. In short, each node of this \textit{binary tree} contains a point. To construct the tree,
let's assume we're at the root. To decide which point gets to be the root, we'll take the \textit{median}
point based on the first axis (in two dimensions that'll just be $x$). Now, all points with $x$ less
than the median go to the left subtree and the rest go to the right subtree. Recurse down to
the next subtree and this time, we'll partition based on the second axis (i.e., $y$). Rinse and
repeat till the tree is fully constructed. Note that the maximum depth of this tree might \textit{not}
be $d$, it could exceed $d$ and if so we just loop back to the first axis.

To resolve a query, let us conduct a depth-first search from the root.
At each node in the tree, let's use it to improve the current score. Next, we need to figure
out which subtree is likely to give us a better answer, and the instinct is to use the current
axis as a basis of comparison. For example, if $x$ of our query point is greater than that of the
current node, let's recurse the right subtree. What's important to realise is that this
is just a \textit{heuristic}, it's not guaranteed that the nearest point is going to be in
the right subtree!

Denoting the gap along this axis as $g = x_{q} - x_{m}$, where $x_{q}$ is the $x$ coordinate
of the query point and $x_{m}$ is the $x$ coordinate of the node, and without loss of generality
assume that $g \geq 0$ (the right subtree case), then if $g \geq c_{\text{best}}$, where
$c_{\text{best}}$ denotes the current minimum distance, there's no point exploring the
left subtree since we'll need to pay the cost of $g$ which already exceeds our best answer.

Empirically, this data structure works great on our points because of the underlying uniform
random distribution. During tree construction, we're almost always able to attain two subtrees
of approximately equal size. As a result, the binary tree is very balanced, and the build stage
consumes $\mathcal{O}(n \log n)$ time. Another phenomenon at low dimensions
is that this pruning heuristic cuts off useless subtrees really well, allowing us to enjoy
point queries in roughly $\mathcal{O}(\log n)$ time, leading to a total time complexity of $\mathcal{O}(n \log n)$.
Of course, things break down at higher dimensions because points become sparse,
and the contribution of each dimension towards the overall distance is tinier. Eventually,
we end up needing to search both subtrees. Fortunately, it just degrades to the brute force,
unlike general grid decomposition that mandates searching all neighbouring hypercubes.
